\section{Introduction}
A numerical decision rule is useful for economic analysis only when its error is small relative to the comparison being made. This requirement is distinct from a small training loss. It is also distinct from outperforming another learning algorithm. When a decision maker can change a preference state, an apparent welfare gain may reflect either the adjustment technology or an inaccurate evaluation of the policy without that technology. When termination is an economic contract, an inaccurate treatment of first exit can dominate both effects.

This paper studies Neural Bellman Operators (NBO): separated policy evaluation and feasible policy improvement, implemented by neural networks and assessed by independent verification. The motivating economy combines consumption, portfolio choice, stochastic risk aversion, and costly preference adjustment. The analysis retains its original running utility, continuous controls, correlated Brownian shocks, and first-exit liquidation covenant. Neither reflecting the state at a training boundary nor altering the running payoff is used to establish the economic results below.

The main numerical result is a certified $.01$-optimal policy in the original continuous economy at the central initial state, for adjustment-cost coefficients $.5$, $2$, and $8$. The respective regret upper bounds are $.00987440$, $.00971727$, and $.00996896$. Each lower endpoint evaluates a feasible policy; each upper endpoint covers all original adaptive controls. The computation does not replace the first-exit contract, restrict the original optimum to time controls, or use a finite action grid to claim a continuous guarantee.

The three original cost certificates can also be turned into a policy rule for an entire cost interval. We construct a library of eighteen certified policies and prove, at the same initial state, a regret bound of $.00999403$ for every real $k\in[.5,8]$. The acceptance test is finite but its guarantee is not restricted to a parameter grid. Convex chords bound the adaptive optimum, and sign-aware expenditure enclosures transfer each feasible policy to all other costs. Exact rational breakpoint tests identify the largest gap between these bounds. This supplies an executable continuation in an economically meaningful parameter without altering the payoff, control opportunities, or stopping contract.

The economic consequences are quantitative. Raising the adjustment-cost coefficient from two to eight lowers optimal welfare by an amount in $[.02318257,.04286880]$. The width of this interval is smaller than its lower endpoint. Access to adjustment at coefficient two raises optimal welfare by an amount in $[.04096859,.05973218]$ relative to the \emph{optimal} no-adjustment economy, rather than relative to a selected no-adjustment rule. Discounted expenditures of the deployed policies are also enclosed independently. The optimal expenditure itself is bounded by valid secants, not identified by relabeling a learned policy as an optimizer.

The computational argument separates policy generation from proof. A learned time-control actor is polished within the feasible output class and evaluated by Gaussian polynomial moments under exact first exit. An affine-preference market-deflator witness then bounds the flexible optimum. Portfolio terms cancel for every adapted original portfolio rule. A supporting source inequality retains flexible adjustment, and a conditional-variance correction pays explicitly for its random marginal value. Outward weighted Taylor boxes enclose the remaining one-Gaussian integrals. Explicit continuation potentials cover artificial verification faces without altering the economic domain.

These constructions improve on the preceding wealth-independent witness, whose $k=2$ regret upper was $.10634$. That witness and its complete proofs remain available; the new certificate is not a relabeling of its evaluation precision as optimality precision. The restricted economy's separate market-deflator bound is retained, together with its independently evaluated feasible policy. It gives an initial restricted-value interval narrower than $.00905$. Combining the two economies' bounds establishes the welfare inequalities above.

A second contribution addresses the inner optimization needed for verification. In a coupled nonconvex control family, all nonconvex coupling depends on the sum of the actions. A one-dimensional branch-and-bound calculation, together with strongly convex scalar dual subproblems, certifies the full continuous action box without a tensor action grid. Executed certificates reach gaps below $10^{-5}$ through 128 action dimensions. At the eight- and sixteen-dimensional learned-critic test states, interval differentiation of the actual network weights also encloses the error in the floating-point critic gradient. This is an implemented bridge between learned parameters and an action certificate, not an assumption that an optimizer's stationary point is globally greedy.

The same nonconvex family is given an absolute accuracy axis. A pointwise running-cost relaxation and removal of control bounds admit a Cole--Hopf representation that retains the original nonquadratic terminal payoff. Orthogonal Gaussian coordinates integrate analytically, leaving a positive scalar integral with an outward enclosure. The resulting number is a lower bound for the unchanged original minimization problem. Relative performance is evaluated separately in a frozen holdout experiment: three dimensions, two training budgets, twelve seed pairs per panel, equal learning-rate search budgets, warmed initialization, randomized method order, and a third, independently specified feedback rule. The pinned author implementation of SOC-MartNet is not modified.

Classical verification, monotone approximation, and primal--dual numerical methods supply important foundations. The stopped It\^o argument is not claimed as a new general verification principle. Nor is alternating neural evaluation and improvement claimed as a new algorithmic idea. The contribution is the constructive combination of original-contract verification, independent economic upper and lower bounds, computable neural-jet allowances, and a non-enumerative global action audit, with each error quantity matched to an executed calculation. Section~\ref{r9:related} compares these objects theorem by theorem with neighboring methods. Recursive utility, sophisticated temporal selves, dynamic games, stochastic-trace identities, counterexamples, and the complete preceding manuscripts are preserved in the supplement rather than being used to enlarge the empirical claims of the present application.


The present revision also tests two previously separate parts of the numerical argument. A new, prospectively frozen comparison expands method-specific tuning under equal nominal training opportunities, implements a predeclared boundary-extension rule, and uses entirely new holdout seeds. Independently reference-solvable tracking, mean-reverting, and ill-conditioned control problems give absolute-accuracy frontiers for implemented held-action actors. A two-error identity separates deployment discretization from actual actor output error. All original-economy results and the complete earlier external holdout remain visible, so neither a stronger comparison protocol nor an easier reference family can silently replace the motivating economic calculation.

\section{The economic environment}\label{r9:economy}
The state is $s=(u,x)$, where $u$ denotes relative risk aversion and $x$ denotes wealth. An admissible action is $a=(c,\theta,p)$ in the compact set
\[
 A=[.05,.8]\times[-.2,.2]\times[-.5,.8].
\]
Controls are progressively measurable and have well-defined stopped state equations. The feedbacks and time controls actually evaluated below are bounded and satisfy this requirement. Before termination,
\begin{align}
 du_t&=\theta_t\,dt+.05\,dZ^u_t,\\
 dx_t&=[(.02+.06p_t)x_t-c_t]dt+.2p_tx_t\,dZ^x_t,
 \qquad d\langle Z^u,Z^x\rangle_t=-.25\,dt.\label{r9:dynamics}
\end{align}
The horizon is $T=1$, discounting is $\rho=.04$, and the open state domain is
$D=(1.2,2.8)\times(.5,2)$. Write $\tau=T\wedge\inf\{r\geq t:S_r\notin D\}$. For $k>0$, running utility and settlement are
\begin{align}
 \ell_k(u,c,\theta)&=\frac{c^{1-u}}{1-u}-\frac{k}{2}\theta^2,\\
 G(u,x)&=-.02(u-2)^2+.1\log x,
 \qquad g(t,u,x)=G(u,x)-8(1-t).\label{r9:payoff}
\end{align}
The policy payoff and optimal value are
\begin{equation}
 J_k^\pi(t,s)=\E^\pi_{t,s}\!\left[\int_t^\tau e^{-\rho(r-t)}\ell_k(S_r,a_r)\,dr
        +e^{-\rho(\tau-t)}g(\tau,S_\tau)\right],
 \qquad V_k=\sup_\pi J_k^\pi.\label{r9:objective}
\end{equation}
The no-adjustment value $V_R$ restricts $\theta$ to zero and retains \emph{all} admissible consumption and portfolio choices. It is independent of $k$. The cost experiment holds the controls, transition law, stopping rule, settlement, utility normalization, and consumption units fixed. A state-dependent transformation of utility would generally define another economy.

The early settlement fee is a liquidation covenant. There is no continuation consumption or wealth injection following exit. At the lower wealth boundary the largest admissible wealth drift is negative. At the upper boundary, $p=0$ eliminates the wealth diffusion and permits an inward drift. Thus neither a silently viable rectangle nor a theorem requiring uniform controlled-boundary regularity can be substituted for this contract. The verification below requires only inequalities for smooth test functions on the trace at which the process is actually stopped.

For a smooth function $v$, put
\begin{align}
 \mathcal L^av={}&\theta v_u+[(.02+.06p)x-c]v_x
   +.00125v_{uu}+.02p^2x^2v_{xx}-.0025pxv_{ux},\\
 R_av={}&v_t+\mathcal L^av+\ell_k-\rho v.\label{r9:generator}
\end{align}
The mixed derivative is retained. Its sign is not inferred solely from the negative shock correlation.

\begin{proposition}[Complete selection at a fixed jet]\label{r9:selection}
For the generator in \eqref{r9:generator}, a global maximizing action can be selected componentwise. Consumption is $\operatorname{proj}_{[.05,.8]}(v_x^{-1/u})$ when $v_x>0$ and is $.8$ otherwise. Adjustment is $\operatorname{proj}_{[-.2,.2]}(v_u/k)$. Let $h_1=.06xv_x-.0025xv_{ux}$ and $h_2=.04x^2v_{xx}$. If $h_2<0$, the portfolio is $\operatorname{proj}_{[-.5,.8]}(-h_1/h_2)$; if $h_2\geq0$, an endpoint maximizing $h_1p+h_2p^2/2$ is selected, with a fixed tie rule.
\end{proposition}
The statement includes convex and zero-curvature portfolio objectives and nonpositive marginal wealth value. It applies to the continuous generator, not to the interpolated, killed-Euler finite-step target. It says nothing about global optimization over shared neural parameters.

\section{Learning and verification are separate operations}\label{r9:algorithm}
For a fixed policy, exact evaluation solves $R_\pi v=0$ with the stopping contract. Exact improvement maximizes the Hamiltonian of that fixed evaluation. Denote these operators by $\mathcal E$ and $\mathcal I$; ideal policy iteration is $\pi_{j+1}=\mathcal I(\mathcal E\pi_j)$. A neural implementation approximates both operators, while an independent evaluator measures the payoff of the policy that will actually be used.

In the raw NBO implementation, critic updates hold the actor fixed. Actor updates hold critic parameters fixed, but retain differentiation through the controlled successor state. A rival's parameters, when present in a game, are also frozen during unilateral improvement. This separation prevents an actor loss from lowering itself by changing the object being evaluated. It does not make the joint network optimization convex. The ideal policy-iteration compactness and limiting arguments are retained in the supplement; no such argument is invoked to certify an observed Adam trajectory.

The inherited finite safeguard is a distinct algorithm. For $\mathcal T_n^aw=r_n^a+\delta P_n^aw$, positive substochastic $P_n^a$, and $\delta=e^{-\rho\Delta}$, define the policy-evaluation residual $e_n=\|v_n-\mathcal T_n^{\pi_n}v_{n+1}\|_\infty$ and the full candidate improvement gap $\eta_n=\|\max_a\mathcal T_n^av_{n+1}-\mathcal T_n^{\pi_n}v_{n+1}\|_\infty$. If the terminal error is at most $b$, backward comparison yields
\begin{equation}
 \|V_{h,n}^*-V_{h,n}^{\pi}\|_\infty
 \leq2\delta^{N-n}b+\sum_{j=n}^{N-1}\delta^{j-n}(2e_j+\eta_j).\label{r9:finite}
\end{equation}
The state-dependent envelope and threshold-correction proof are reproduced in the supplement. A threshold can reduce the fraction of decisions changed while leaving the cost of a complete candidate audit essentially unchanged. In particular, the retained 125-action tensor audit is not the non-enumerative certifier of Section~\ref{r9:action}. Its floating-point envelope is not relabeled as an outward continuous-time bound.

\subsection{A quantitative bridge from a fitted network}
Let $C_\rho(h)=(1-e^{-\rho h})/\rho$, with $C_0(h)=h$.
\begin{assumption}[Admissible witnesses]\label{r9:witness}
Witnesses are bounded, continuous on the stopped trace, and $C^{1,2}$ in the interior of each time slab. Localized It\^o identities pass to the stopping time by uniform integrability. At slab interfaces an upper witness has nonpositive forward jumps and a lower witness has nonnegative forward jumps. The implemented policy has a well-defined stopped state equation.
\end{assumption}
\begin{theorem}[Stopped one-sided comparison]\label{r9:pair}
Suppose Assumption~\ref{r9:witness} holds, $U\geq g-b_U$ and $L\leq g+b_L$ on the stopped trace, $\sup_aR_aU\leq\epsilon_U$, and $R_\pi L\geq-\epsilon_L$, where all four allowances are nonnegative. Then
\begin{align}
 J_k^\pi&\geq L-b_L-C_\rho(T-t)\epsilon_L,\qquad
 V_k\leq U+b_U+C_\rho(T-t)\epsilon_U,\\
 0\leq V_k-J_k^\pi&\leq U-L+b_U+b_L+C_\rho(T-t)(\epsilon_U+\epsilon_L).\label{r9:pairbound}
\end{align}
The inequalities do not require an attained optimizer or a smooth optimal value.
\end{theorem}

\begin{theorem}[Cover-and-oracle certificate for neural output]\label{r9:bridge}
Let a continuous neural witness $v$ and a feasible neural policy $\pi$ satisfy Assumption~\ref{r9:witness}, with stopped-trace error at most $b$. Let cells cover the full audited time--state domain and have representatives $z_j$. Suppose independent calculation proves, throughout cell $j$,
\[
 |R_\pi v(z)-R_\pi v(z_j)|\leq\omega_j,
 \quad |\Gamma(z)-\Gamma(z_j)|\leq\nu_j,
 \qquad \Gamma=\sup_aR_av-R_\pi v.
\]
Suppose the representative evaluation intervals imply $|R_\pi v(z_j)|\leq e_j$ and a global action oracle implies $\Gamma(z_j)\leq q_j$, including jet, inner-solve, and arithmetic allowances. Set $e=\max_j(e_j+\omega_j)$ and $q=\max_j(q_j+\nu_j)$. Then
\begin{equation}
 0\leq V_k-J_k^\pi\leq 2b+C_\rho(T-t)(2e+q).\label{r9:bridgebound}
\end{equation}
In particular, certified derivative bounds $M_j,N_j$ and cell radii $h_j$ permit $\omega_j=M_jh_j$ and $\nu_j=N_jh_j$. Any sequence of learned outputs for which these verified allowances vanish has vanishing regret, irrespective of whether its parameters converge.
\end{theorem}
This is a quantitative a posteriori connection to network errors, not a claim that an empirical mean-square loss controls a supremum. Interval differentiation or weight-based derivative bounds can supply the cell moduli. They may be conservative, and covering a high-dimensional state domain remains costly. Section~\ref{r9:action} executes the jet-to-action part for actual learned critics. The economic calculation instead fits an independent polynomial witness; it does not falsely claim to have certified a whole-state neural critic merely because its action audit passed at selected states.

\section{A constructive certificate for the flexible economy}\label{r10:flexible}
The principal numerical target is a regret bound of $.01$ for the original problem \eqref{r9:objective}, not for a finite approximation. The certificate below attains this target at $(0,2,1.25)$ for each of the three adjustment prices. It uses a feasible policy to bound value from below and a separate, affine-preference market-deflator construction to bound the optimum from above. The latter retains arbitrary adapted consumption, portfolio choice, and preference adjustment. Neither a state lattice nor an optimal-value label enters its construction.

\subsection{An affine source that preserves the original opportunities}
Set $\phi(u)=-.02(u-2)^2$, $h=1-t$, and let $Y$ satisfy \eqref{r9:deflator}, with its initial level $y_0>0$ now chosen as a witness parameter. Define
\begin{align}
 S_0(u,y)&=\Psi(u,y)+.01y-.004\log(.5),\\
 \beta(u,y)&=\partial_u\Psi(u,y),\\
 m(t)&=2+\int_0^t\bar\theta_s\,ds,\qquad 0\leq\bar\theta_s\leq.2,\label{r10:source}
\end{align}
where $\Psi$ is the original consumption conjugate in \eqref{r9:dualsource}. The reference drift $\bar\theta$ is deterministic. It is a parameter of an upper witness, not a restriction on the decision maker's adjustment policy. Write $\widehat y=\operatorname{proj}_{[.2,12]}y$ and
\[
 Q(u)=\phi(m(1))+\phi'(m(1))(u-m(1)),\qquad b_T=\phi'(m(1)).
\]
The verification rectangle is $u\in[1.5,2.5]$, $y\in[.2,12]$. On it, the following supporting inequality holds:
\begin{equation}
 S_0(u,y)\leq S_0(m,y)+\beta(m,y)(u-m),\quad 2\leq m\leq2.2.\label{r10:plane}
\end{equation}
It is a one-sided inequality on the whole rectangle, not a fitted tangent at sampled states. Its proof combines derivative identities and interval inequalities with explicit margins; see Appendix~\ref{r10:proof:plane}. Global concavity of $S_0$ outside this rectangle is not assumed.

Put $\widehat S_r(u,y)=S_0(m(r),\widehat y)+\beta(m(r),\widehat y)(u-m(r))$. Consider the full-line auxiliary control problem
\begin{align}
 A_k(t,u,y)=\sup_{|\theta|\leq.2}\E_{t,u,y}\bigg[&\int_t^1 e^{-\rho(r-t)}
 \{\widehat S_r(U_r,Y_r)-k\theta_r^2/2\}dr\notag\\
 &+e^{-\rho(1-t)}Q(U_1)\bigg].\label{r10:auxiliary}
\end{align}
Here $dU_r=\theta_rdr+.05dZ^u_r$, and the auxiliary control remains adapted. The affine source makes its marginal value of $u$ an explicitly defined conditional expectation, which can be bounded without solving another neural control problem. Put
\begin{align}
 Z_u(h,u)&=8h\{e^{-120(u-1.5)+42h}+e^{-120(2.5-u)+42h}\},\\
 Z_y(h,y)&=8h\{e^{-22\log(y/.2)+22.33h}+e^{-22\log(12/y)+22.33h}\}.\label{r10:localization}
\end{align}

\begin{theorem}[Affine-preference dual and a computable initial upper bound]\label{r10:dual}
For $k\geq.5$, $y_0\in[1.2,2.4]$, and any deterministic $\bar\theta\in[0,.2]$, the original flexible optimum satisfies
\begin{equation}
 V_k(0,2,1.25)\leq .75y_0+.1\log(.5)+A_k(0,2,y_0)+Z_u(1,2)+Z_y(1,y_0).\label{r10:dualbound}
\end{equation}
To evaluate the right side, put $\beta_t=\beta(m(t),\widehat Y_t)$ and
\begin{align}
 b(s)&=\int_s^1e^{-\rho(t-s)}\E\beta_t\,dt+e^{-\rho(1-s)}b_T,\\
 g_k(q,a)&=H_k(q)-qa+ka^2/2,\\
 C&=.00375\int_0^1te^{-\rho t}\E\left[\partial_l\beta(m(t),\widehat Y_t)\right]dt,\label{r10:quantities}
\end{align}
where the $l$ derivative includes the clipping map and $H_k(q)=\max_{|a|\leq.2}(qa-ka^2/2)$. Then
\begin{align}
 A_k(0,2,y_0)\leq{}&e^{-\rho}\phi(m(1))+
 \int_0^1e^{-\rho t}\left\{\E S_0(m(t),\widehat Y_t)-k\bar\theta_t^2/2\right\}dt+C\notag\\
 &+\int_0^1e^{-\rho s}g_k(b(s),\bar\theta_s)ds+\frac{.09D^2}{24k},\label{r10:computable}
\end{align}
with the explicit bound
\begin{equation}
 D^2=\frac{e^{aM+.045a^2}}{16}\{(M+.09a)^2+.09\},\qquad
 a=12/11,\quad M=\log y_0.\label{r10:variance}
\end{equation}
All expectations in \eqref{r10:quantities}--\eqref{r10:computable} are one-Gaussian integrals. They involve no unknown optimal value or policy.
\end{theorem}

The cancellation of the portfolio is the same accounting identity used in the restricted dual: $d(Yx)=(.04Yx-Yc)dt+$ a local martingale. The new step is to retain flexible adjustment while replacing its source by a supporting affine function. The variance allowance in \eqref{r10:computable} pays explicitly for the random conditional marginal value of adjustment. Replacing that marginal value by its unconditional mean without this allowance would generally produce an invalid upper bound. The covariance term $C$ also matters: it retains the original correlation between preference and market shocks and is nonpositive here. It is evaluated, not silently discarded or confused with a change of measure.

The source bound $S_0\geq-7.1$ and $|\beta|\leq1.1$ imply $A_k-Q\geq-7.925272h$ on the verification rectangle, by choosing zero adjustment in the auxiliary problem. Hence the factors $8h$ in \eqref{r10:localization} pay for continuation at artificial faces. On actual wealth exits, the uncorrected affine witness already dominates the liquidation settlement. The global fallback $V_k\leq G$ covers continuation after an artificial exit. These facts connect \eqref{r10:computable} to the original stopped economy rather than to a surrogate full-line economy.

\subsection{Policy generation, verified refinement, and the stopping rule}\label{r10:algorithm}
The lower-bound policies start from the retained one-input, twelve-unit neural time-control actors. A constrained L-BFGS-B step polishes their sixteen consumption and adjustment outputs using the original payoff and the same exact budget projection. This is an explicit enlargement from a fixed network's image to the feasible time-control outputs; it is not described as additional Adam iterations or as a theorem that the neural optimizer converges. Consumption and adjustment constants are rounded inward relative to the \emph{exact decimal} action bounds, and feasibility is checked using their exact binary rational values. The finite-node training objective is a proposal mechanism, not an asserted untruncated Gaussian expectation of singular CRRA utility. The independent polynomial-moment evaluator in Appendix~\ref{r9:proof:policy} then encloses each deployed continuous-time policy under exact first preference exit. The full network histories, original outputs, polished outputs, and output-optimization records remain separate.

The dual reference is fitted with Gaussian quadrature, but its fitted objective is never used as a bound. Verification uses $t=v^2$, subdivides each of the sixteen time slabs, and covers the standard Gaussian coordinate on $[-6,6]$. Constant and linear Taylor terms are integrated with interval-enclosed weighted moments. Hessian ranges enclose the entire remainder, including boxes crossing the consumption constraint. Gaussian tails, the terminal tangent, the conditional-variance allowance, localization, and outward arithmetic are separately included. The small one-dimensional interval tests establishing \eqref{r10:plane} precede every certificate execution.

The stopping rule is simple. Given a policy lower bound $L$ and an upper bound $U$ from Theorem~\ref{r10:dual}, return the policy with tolerance $\varepsilon$ only when $U-L\leq\varepsilon$. Otherwise refine the quadrature or return the explicit unachieved target at the declared work limit. Soundness does not require optimizer convergence. Nor does soundness imply that an arbitrary tolerance can be reached with a fixed affine relaxation. For the three declared $.01$ instances, the finite executions below supply the termination certificate.

\begin{table}[t]\centering\small
\caption{Original flexible-economy certificates after output polishing}\label{r10:tab:original}
\begin{tabular}{rrrr}
\toprule
$k$ & Feasible policy payoff & Optimal-value upper & Regret upper \\
\midrule
0.5 & $[-1.27267126,-1.27265474]$ & -1.26279686 & 0.00987440 \\
2 & $[-1.29502451,-1.29501124]$ & -1.28530723 & 0.00971727 \\
8 & $[-1.32817604,-1.32816919]$ & -1.31820708 & 0.00996896 \\
\bottomrule
\end{tabular}

\par\smallskip\footnotesize All three regret bounds are below $.01$. Each upper bound covers the full original continuous control set and stopping contract. Each lower bound evaluates the specified feasible time-control policy. Endpoints are rounded outward; gaps use the unrounded endpoints. These are central-state certificates, not uniform claims over every initial condition.
\end{table}

\begin{table}[t]\centering\small
\caption{Verified work needed to cross the absolute-error target}\label{r10:tab:frontier}
\begin{tabular}{rrrrr}
\toprule
$k$ & Source boxes & Regret upper & Verification seconds & $.01$ met \\
\midrule
0.5 & 4,096 & 0.01027899 & 5.46 & No \\
0.5 & 16,384 & 0.00995838 & 11.73 & Yes \\
0.5 & 65,536 & 0.00987440 & 27.54 & Yes \\
2 & 4,096 & 0.01013162 & 5.25 & No \\
2 & 16,384 & 0.00980424 & 11.30 & Yes \\
2 & 65,536 & 0.00971727 & 27.06 & Yes \\
8 & 4,096 & 0.01037189 & 5.29 & No \\
8 & 16,384 & 0.01005494 & 11.48 & No \\
8 & 65,536 & 0.00996896 & 26.79 & Yes \\
\bottomrule
\end{tabular}

\par\smallskip\footnotesize The policy is independently evaluated once at the finest stated precision. Time includes that evaluation, dual preparation, and cumulative dual verification through each row, on one CPU thread. It excludes inherited network training, witness fitting, imports, and setup. This is an absolute-accuracy \emph{verification} frontier, not an all-in training comparison or an external matched-accuracy superiority claim. Raw stage clocks are retained.
\end{table}

\begin{table}[t]\centering\scriptsize
\caption{Separate upper-bound allowances at $k=2$}\label{r10:tab:allowances}
\begin{tabular}{rrrrr}
\toprule
Boxes & Source remainder & Control gap & Variance allowance & Localization \\
\midrule
4,096 & 0.000339131 & 0.000063959 & 0.000081080 & 0.000000247 \\
16,384 & 0.000082675 & 0.000041416 & 0.000081080 & 0.000000247 \\
65,536 & 0.000020398 & 0.000032100 & 0.000081080 & 0.000000247 \\
\bottomrule
\end{tabular}

\par\smallskip\footnotesize The source remainder is the Taylor allowance on the covered Gaussian region. Analytic Gaussian tails, interval weight errors, terminal evaluation, and the covariance interval are included in the full records. The variance allowance does not decrease merely by subdividing quadrature boxes. No row is obtained by treating a mesh difference as a rigorous error estimate.
\end{table}

The finest upper bounds are approximately $-1.26279687$, $-1.28530724$, and $-1.31820708$. Combined with the feasible-policy lower endpoints, they give regret upper bounds $.00987440$, $.00971727$, and $.00996896$ after outward display rounding. The central target is attained without a finite candidate-action grid, an Euler-error assumption, or an assumption that the policy is optimal within its neural class. The preceding wealth-independent witness and its larger gaps are retained in Appendix~\ref{r10:retained}, where their loss of wealth information can be compared with the new construction.

\subsection{A quantitative bridge from implemented controls to the certificate}
The bound applies to the stored deployment constants themselves. A further perturbation result quantifies its robustness to control evaluation error, rather than inferring robustness from a small training loss.
\begin{proposition}[Feasible-output error allowance]\label{r10:output}
Consider two deterministic time-control policies with $p=0$, $c,\widetilde c\in[.7,.8]$, $\theta,\widetilde\theta\in[0,.2]$, and strictly feasible budgets as in \eqref{r9:budgetprojection}. Let
\[
 \delta_c=\int_0^1|c_t-\widetilde c_t|dt,\qquad
 \delta_\theta=\int_0^1|\theta_t-\widetilde\theta_t|dt.
\]
Their original stopped payoffs satisfy
\begin{equation}
 |J_k(c,\theta)-J_k(\widetilde c,\widetilde\theta)|
 \leq(3+.2e^{.02})\delta_c+(25+.2k+.032)\delta_\theta
           +8(14.1+.02k)e^{-72}.\label{r10:outputbound}
\end{equation}
Consequently, a validated output error bounded by the right side is an additive allowance in the returned policy's regret certificate. Feasibility must still be checked separately; small unconstrained control error alone does not rule out early wealth exit.
\end{proposition}
For the stored policies the output allowance is zero: their binary64 constants are treated as exact reals by the evaluator. Proposition~\ref{r10:output} instead gives an explicit interface for a budget-preserving implementation with perturbed outputs. Together with the neural-jet bound of Theorem~\ref{r9:bridge}, it connects observable implementation errors to certified quantities without attributing an unproved global convergence rate to Adam, RMSprop, or L-BFGS-B.


\section{Certifying the value of preference adjustment}\label{r9:economics}
For each feasible policy write $J_k(\pi)=A(\pi)-kB(\pi)$, where
\[
 B(\pi)=\E^\pi\int_t^\tau e^{-\rho(r-t)}\theta_r^2/2\,dr,
 \qquad 0\leq B(\pi)\leq .02C_\rho(T-t).
\]
The policy set is common to every $k$.
\begin{theorem}[Accuracy-calibrated cost comparisons]\label{r9:cost}
The optimal value is convex and nonincreasing in $k$. If optimizers exist, their adjustment budgets are nonincreasing across distinct prices. More generally, let $\pi_i$ be $\varepsilon_i$-optimal at $k_i$, with $k_2>k_1$. Then
\begin{equation}
 B(\pi_2)-B(\pi_1)\leq\frac{\varepsilon_1+\varepsilon_2}{k_2-k_1}.\label{r9:reversal}
\end{equation}
If $V(k)\in[\underline V(k),\overline V(k)]$, the budget of an optimizer at $k$, for $k_-<k<k_+$, lies in the primitive budget range intersected with
\begin{equation}
 \left[\frac{\underline V(k)-\overline V(k_+)}{k_+-k},
       \frac{\overline V(k_-)-\underline V(k)}{k-k_-}\right].\label{r9:secant}
\end{equation}
\end{theorem}
The theorem concerns discounted adjustment expenditures, not a pointwise ordering of controls. A re-solved $V_u/k$ is not a comparative-static proof. The new full-value enclosures in Table~\ref{r10:tab:original} now resolve strict welfare losses across the cost grid. At $k=2$, \eqref{r9:secant} also gives the optimal-budget enclosure $[.00386376,.01960529]$, conditional on existence of an optimizer. Its positive lower endpoint is informative, while its upper endpoint is still the primitive bound. The tightly enclosed expenditures in Table~\ref{r9:tab:economics} describe the specified learned-and-polished policies, not an identified optimal budget.

\subsection{A market-deflator witness for the restricted optimum}\label{r9:dual}
A tighter comparison becomes possible by exploiting the restricted economy's structure. Introduce the auxiliary process
\begin{equation}
 dY_t=.02Y_t\,dt-.3Y_t\,dZ^x_t,\qquad Y_0=1.55.\label{r9:deflator}
\end{equation}
It is a test-process constructed from the original shocks. It does not restrict portfolio information, change wealth, or add a traded asset. For
\[
 W_0(u,x,y)=y(x-.5)+.1\log(.5)-.02(u-2)^2,
\]
the $p$ terms in the generator of $Yx$ cancel exactly. With $\theta=0$, define
\begin{align}
 \Psi(u,y)&=\max_{.05\leq c\leq.8}\left\{\frac{c^{1-u}}{1-u}-yc\right\},\\
 S(u,y)&=\Psi(u,y)+.01y-.00005-.004\log(.5)+.0008(u-2)^2.\label{r9:dualsource}
\end{align}
On $u\in[1.5,2.6]$, $y\in[.2,12]$, one has $-8<S<0$. Extend $S$ to a bounded function $\widetilde S\in[-8,0]$ by clipping its two arguments at these verification faces. Let
\begin{equation}
 F_R(h,u,y)=\E_{u,y}\int_0^h e^{-.04r}\widetilde S(U_r,Y_r)\,dr,\label{r9:restrictedpotential}
\end{equation}
where $dU_r=.05\,dZ^u_r$ and $Y$ follows \eqref{r9:deflator}. Thus $-8h\leq F_R\leq0$.

\begin{theorem}[A restricted-economy upper bound]\label{r9:restricted}
Let $Z_R$ be the sum of the following nonnegative localization potentials:
\begin{align}
 Z_{R,u}&=8h\{e^{-200(u-1.5)+50h}+e^{-240(2.6-u)+72h}\},\\
 Z_{R,y}&=8h\{e^{-22\log(y/.2)+22.33h}
                         +e^{-22\log(12/y)+22.33h}\}.
\end{align}
Then
\begin{equation}
 V_R(0,2,1.25)\leq W_0(2,1.25,1.55)+F_R(1,2,1.55)+Z_R(1,2,1.55).
 \label{r9:restrictedupper}
\end{equation}
The upper bound covers all adaptive original consumption and portfolio controls with $\theta=0$. Consequently, for any feasible flexible policy with certified payoff lower bound $\underline J_k$ and any upper bound $\overline V_R$ from \eqref{r9:restrictedupper},
\begin{equation}
 V_k(0,2,1.25)-V_R(0,2,1.25)\geq\max\{0,\underline J_k-\overline V_R\}.\label{r9:accesslower}
\end{equation}
\end{theorem}

\begin{table}[t]
\centering\small
\caption{Refining the no-adjustment dual witness}\label{r9:tab:dual}
\begin{tabular}{r r r r}
\toprule
Weighted boxes & Dual witness-value enclosure & Width & Seconds \\
\midrule
16384 & $[-1.33876765,-1.33309541]$ & 0.00567 & 2.67 \\
131072 & $[-1.33676857,-1.33544680]$ & 0.00132 & 15.13 \\
1048576 & $[-1.33631272,-1.33599309]$ & 0.00032 & 110.24 \\
\bottomrule
\end{tabular}

\par\smallskip\footnotesize These intervals enclose the \emph{dual witness value}, not the restricted optimum. Only their upper endpoints are used as upper bounds for that optimum. A separately evaluated restricted policy gives the primal lower bound. Original stopping and all adaptive consumption and portfolio opportunities are retained.
\end{table}

The expectation in \eqref{r9:restrictedpotential} has an explicit two-Gaussian representation. We transform time by $r=v^2$, evaluate constant and linear Taylor terms with interval-enclosed weighted moments, and bound the Hessian remainder on each box. Exponential Taylor remainders, Gaussian tails, and localization are included. This produces the executed refinement in Table~\ref{r9:tab:dual}; it does not appeal to a mesh difference as an error estimate. The final restricted-optimum upper bound is $-1.33599309$ after outward display rounding. The evaluated restricted policy has payoff in $[-1.34503942,-1.34503585]$, so the restricted problem itself has an initial optimal-value enclosure narrower than $.00905$.

\begin{table}[t]
\centering\scriptsize
\caption{Certified current-policy expenditures and optimal value of access}\label{r9:tab:economics}
\begin{tabular}{rrrr}
\toprule
$k$ & Policy adjustment budget & Optimal access-value interval & Regret target \\
\midrule
0.5 & $[0.01757696,0.01757697]$ & $[0.06332184,0.08224255]$ & $<.01$ \\
2 & $[0.01241866,0.01241867]$ & $[0.04096859,0.05973218]$ & $<.01$ \\
8 & $[0.00192681,0.00192682]$ & $[0.00781706,0.02683234]$ & $<.01$ \\
\bottomrule
\end{tabular}

\par\smallskip\footnotesize The expenditure column evaluates the displayed learned policy. The access columns bound the difference between the two \emph{optimal} values. The lower endpoint pairs a feasible flexible policy with the restricted dual upper; the upper endpoint pairs the flexible upper with a feasible restricted policy. All endpoints are rounded outward from full-precision records. The cost-two access interval has width below its entire positive gain; at cost eight, the interval still resolves the sign but not a comparably precise magnitude.
\end{table}

The welfare lower bounds remain strictly positive at all three prices. At $k=2$ the optimal access gain lies in $[.04096859,.05973218]$; its width is smaller than the lower endpoint. Table~\ref{r10:tab:cost-effects} separately compares the flexible optima at different adjustment prices. The $2\to8$ and $.5\to8$ effects exceed their entire interval widths. The $.5\to2$ comparison resolves the sign but has a larger relative uncertainty. No comparison assumes that a displayed policy is an optimizer or interprets the value-ordering theorem as a pointwise ordering of adjustment controls.

\begin{table}[t]\centering\small
\caption{Accuracy-calibrated optimal welfare losses from higher adjustment costs}\label{r10:tab:cost-effects}
\begin{tabular}{rrrr}
\toprule
Cost change & Optimal welfare loss & Interval width & Width below loss \\
\midrule
$0.5\to2$ & $[0.01263598,0.03222764]$ & 0.01959166 & Sign resolved \\
$2\to8$ & $[0.02318257,0.04286880]$ & 0.01968623 & Yes \\
$0.5\to8$ & $[0.04553582,0.06537918]$ & 0.01984335 & Yes \\
\bottomrule
\end{tabular}

\par\smallskip\footnotesize Each interval encloses $V_{k_1}-V_{k_2}$ in the unchanged original economy. Bounds use independent value endpoints, not fitted derivatives in $k$. The last column compares the width with the lower endpoint of the effect interval.
\end{table}

\section{A certified policy rule on a continuous cost interval}\label{r12:continuation}
The preceding value and expenditure enclosures can answer a stronger economic question than evaluation at three selected costs: can one supply an approximately optimal policy for \emph{every} adjustment price in an interval? Repeating training on a fine parameter grid is not enough. An unexamined price may lie between grid points, and a policy fitted at one price need not remain approximately optimal at another. The affine dependence of a fixed policy's payoff on the cost coefficient makes a rigorous continuation possible.

Fix $(t,u,x)=(0,2,1.25)$. For ordered prices $k_0<\cdots<k_m$, let $\pi_j$ be a frozen feasible policy and suppose independently computed quantities satisfy
\begin{equation}
 L_j\leq J_{k_j}(\pi_j)\leq V(k_j)\leq U_j,
 \qquad b_j^-\leq B(\pi_j)\leq b_j^+.
 \label{r12:node}
\end{equation}
Here $B$ includes the actual random first exit. It is not an undiscounted or unstopped training proxy. Define
\begin{align}
 \ell_j(k)&=L_j-(k-k_j)
 \begin{cases}b_j^+,&k\geq k_j,\\ b_j^-,&k<k_j,\end{cases}
 &\ell(k)&=\max_{0\leq j\leq m}\ell_j(k),\label{r12:lower}\\
 u(k)&=\frac{k_{i+1}-k}{k_{i+1}-k_i}U_i+
        \frac{k-k_i}{k_{i+1}-k_i}U_{i+1},
 && k\in[k_i,k_{i+1}].\label{r12:chord}
\end{align}
Using the upper expenditure endpoint on \emph{both} sides of a node would be wrong: a lower cost rewards the policy's expenditure, so its lower endpoint must be used on the left.

\begin{theorem}[Certified parameter continuation]\label{r12:envelope}
Suppose the policy set, state dynamics, and stopping contract do not depend on $k$, and \eqref{r12:node} holds. Select the lowest-index maximizer $j(k)$ of \eqref{r12:lower}. For every real $k\in[k_0,k_m]$,
\begin{equation}
 \ell(k)\leq J_k(\pi_{j(k)})\leq V(k)\leq u(k),
 \qquad 0\leq V(k)-J_k(\pi_{j(k)})\leq u(k)-\ell(k).
 \label{r12:uniform}
\end{equation}
On each node interval the maximum of $u-\ell$ occurs at an endpoint or a breakpoint of the upper envelope of the affine functions $\ell_j$. If the inputs are rational, these candidate points and gaps admit an exact finite computation. Strictly bounding all these gaps by $\varepsilon$ certifies \eqref{r12:uniform} on the entire continuous price interval, without a parameter-grid approximation error.
\end{theorem}
The proof in Appendix~\ref{r12:proof} uses the affine fixed-policy payoff and convexity of the optimal value. It does not require differentiability of $V$, an attained optimal control, or a convergence theorem for the network optimizer. This use of convexity is distinct from assuming convexity of the shared-parameter training objective.

\subsection{Construction and stopping rule}
The starting nodes are the three independently certified prices $.5$, $2$, and $8$. A disclosed pilot adds $4$. Thereafter a failed node interval is bisected, in increasing-price order; each new price receives its own feasible actor, exact-first-exit payoff enclosure, expenditure enclosure, and affine-preference dual upper bound. Training uses the nearest inherited actor only as a warm start. Output polishing and dual fitting use the preceding deterministic routines with explicitly audited input/output substitutions. They do not use an optimal-value label. The original source files are unchanged.

All final policy evaluations use 16,384 outward time rectangles, and each optimal-value upper uses 65,536 weighted source boxes. Artificial verification faces carry the same explicit localization allowances as in Theorem~\ref{r10:dual}; none becomes an economic stopping boundary. Each continuous-price test then converts the binary64 certificate endpoints to their exact rational values. A monotone upper-hull algorithm constructs all relevant line intersections. No numerical mesh, random evaluation price, interpolation residual estimate, or optimization tolerance replaces this last test.

\begin{table}[t]\centering\small
\caption{Constructive continuation to the uniform economic tolerance}\label{r12:frontier}
\begin{tabular}{rrrr}\toprule
Certified prices & Uniform regret upper & Task-seconds & Target met \\ \midrule
3 & 0.023989184 & 0.00 & No \\
4 & 0.014155490 & 20.38 & No \\
7 & 0.011507236 & 85.20 & No \\
13 & 0.010387647 & 219.98 & No \\
17 & 0.010002257 & 304.81 & No \\
18 & 0.009994025 & 325.31 & Yes \\
\bottomrule\end{tabular}

\par\smallskip\footnotesize The maximum covers every real price in $[.5,8]$, at the fixed initial state. Task-seconds sum successful new price-case durations, including polishing and verification; parallel tasks overlap. They exclude the inherited three certificates and interrupted attempts and are not end-to-end CPU time. All interruption records are preserved separately. The acceptance criterion is the exact rational gap, not the displayed runtime.
\end{table}

The final library contains eighteen nodes and fifteen newly certified policies. It yields
\begin{equation}
 \sup_{k\in[.5,8]}\{V(k)-J_k(\pi_{j(k)})\}
 \leq .009994024846927508 < .01.
 \label{r12:executed}
\end{equation}
The displayed upper is rounded outward from the exact rational maximum. The largest gap occurs at an envelope intersection near $k=7.7422$, not at a training price. Checking only the eighteen training prices would therefore miss the relevant acceptance calculation. This is also a reason to retain all failed refinement stages rather than show only the last table row.

\begin{proposition}[Finite refinement under local slack]\label{r12:termination}
Suppose each newly supplied node certificate satisfies $U_j-L_j\leq\varepsilon-\eta$ for a common $\eta>0$, and $0\leq b_j^-\leq b_j^+\leq\bar B$. On an interval of length $h$,
\begin{equation}
 \max(u-\ell)\leq\max\{U_i-L_i,U_{i+1}-L_{i+1}\}+\bar B h/4.
 \label{r12:slack}
\end{equation}
Consequently, exact bisection terminates after finitely many refinements when all node intervals have length less than $4\eta/\bar B$.
\end{proposition}
The proposition quantifies the continuation step, not the ability of an arbitrary optimizer to supply the stipulated node slack. A fixed witness class can have a structural accuracy floor. The executed result \eqref{r12:executed} does not rely on an unproved uniform-slack claim: its actual finite library passes the exact test. The tighter upper-hull test, rather than the conservative bound \eqref{r12:slack}, is used for acceptance.

\subsection{Economic interpretation between certified prices}
The library delivers a policy for a queried cost without retraining. It selects one frozen sixteen-slab control rule for the whole horizon and applies the original liquidation contract if exit occurs. It does not switch policies according to an unverified neural value between states. The guarantee is against the full adaptive optimum, not merely against the best policy in the library.

For prices $p<q$, the continuation bounds imply
\begin{equation}
 V(p)-V(q)\in[\ell(p)-u(q),\,u(p)-\ell(q)]\cap[0,\bar B(q-p)].
 \label{r12:effects}
\end{equation}
Similarly, the value of access lies in $[\ell(k)-\overline V_R,u(k)-\underline V_R]$. The original certified access gain at $k=8$ and monotonicity imply a positive access gain throughout $[.5,8]$. This assertion resolves the sign, not a small relative uncertainty for every price. The sharper effect-to-width comparisons at $k=2$ and for the $2\to8$ experiment remain those already reported above.

If an optimizer exists at an interior price $k$, all available secants may be intersected to obtain
\begin{align}
 B(\pi_k^*)&\geq\max\left\{0,\sup_{q>k}\frac{\ell(k)-u(q)}{q-k}\right\},\label{r12:budgetlower}\\
 B(\pi_k^*)&\leq\min\left\{\bar B,\inf_{p<k}\frac{u(p)-\ell(k)}{k-p}\right\}.
 \label{r12:budgetupper}
\end{align}
Using a finite subset of these secants remains valid and avoids numerical differentiation of a possibly nonsmooth value function. The deployed policies' independently enclosed expenditures remain separate objects; no optimal adjustment budget is inferred from them without the requisite optimal-value comparison.

The new object is the executable parameter-uniform policy certificate and its exact stopping rule. Neither convexity of a supremum of affine functions nor classical policy verification is claimed as new. In combination with the original stopped dual, the construction closes an additional economically meaningful computational chain: a continuum of cost counterfactuals, an implementable policy at every cost, and an explicit absolute error smaller than the declared tolerance. It does not convert the separate external fixed-budget experiments into a universal performance claim.


\section{Global action certification without a tensor grid}\label{r9:action}
Consider the following minimization problem, which is also used for the external comparison:
\begin{align}
 dX_t&=2a_t\,dt+\sqrt{.32}\,dW_t,\qquad X_0=0,
       \qquad a_t\in[-1,1]^d,\\
 f_d(a)&=\frac1d\sum_{i=1}^d\{a_i^2+.1\sin^2a_i\}
                 +.3\sin\!\left(\sum_i a_i\right),\\
 q_d(x)&=\frac1d\sum_{i=1}^d(x_i-.5)^2
                 +.2\cos\!\left(d^{-1/2}\sum_i x_i\right),\label{r9:nonconvex}\\
 V_d^*&=\inf_a\E\left[\int_0^1 f_d(a_t)\,dt+q_d(X_1)\right].\notag
\end{align}
The objective has genuinely coupled nonconvex actions. No analytic optimal-action or optimal-value labels are supplied to training. At a fixed critic gradient $p$, the action-dependent Hamiltonian is
\[
 H_p(a)=\sum_i f_i(a_i)+.3\sin z,\qquad
 f_i(s)=\frac{s^2+.1\sin^2s}{d}+2p_is,\quad z=\sum_i a_i.
\]
Each $f_i$ is strongly convex with modulus at least $m=1.8/d$.

\begin{theorem}[Scalar global partition and convex dual bounds]\label{r9:actionthm}
For any multiplier $\lambda$ and $z\in[-d,d]$, put
\[
 D(\lambda)=\sum_i\min_{|s|\leq1}\{f_i(s)-\lambda s\}.
\]
On a scalar interval $z\in[z_0-r,z_0+r]$, a lower bound for the full Hamiltonian is
\begin{equation}
 D(\lambda)+\lambda z_0+.3\sin z_0
       -|\lambda+.3\cos z_0|r-.15r^2.\label{r9:actionlower}
\end{equation}
Any feasible action supplies an upper bound. Each inner convex minimum can itself be bounded below by $q(s_0)-e(s_0)^2/(2m)$, where $q=f_i-\lambda s$ and $e$ is the derivative magnitude in a feasible descent direction, or the full derivative magnitude at an interior point. Outward versions of these bounds, followed by scalar subdivision and pruning, give a global gap on $[-1,1]^d$.

If a stored gradient $\widehat p$ differs from the true neural gradient by componentwise bounds $\delta_i$, the feasible action's true-gradient gap is at most its stored-gradient gap plus $4\sum_i\delta_i$.
\end{theorem}
Floating-point root finding only proposes multipliers and actions. It is not trusted to solve the convex subproblems exactly. Even an inaccurate proposal remains valid after the strong-convexity allowance. This structure removes tensor enumeration; it does not establish a dimension-free complexity theorem or solve a generic nonseparable nonconvex Hamiltonian.

\begin{table}[t]
\centering\small
\caption{Executed continuous-action certificates}\label{r9:tab:action}
\begin{tabular}{r r r r r r}
\toprule
Action dim. & Tolerance & Verified gap & Scalar boxes & Convex subproblems & Seconds \\
\midrule
8 & 1e-03 & 0.000258 & 33 & 264 & 0.299 \\
8 & 1e-05 & 6.72e-06 & 47 & 376 & 0.493 \\
16 & 1e-03 & 0.000931 & 37 & 592 & 0.476 \\
16 & 1e-05 & 4.42e-06 & 57 & 912 & 0.751 \\
32 & 1e-03 & 0.00068 & 85 & 2720 & 1.498 \\
32 & 1e-05 & 3.52e-06 & 103 & 3296 & 1.991 \\
64 & 1e-03 & 0.000647 & 163 & 10432 & 6.041 \\
64 & 1e-05 & 3.06e-06 & 185 & 11840 & 5.955 \\
128 & 1e-03 & 0.000948 & 317 & 40576 & 19.229 \\
128 & 1e-05 & 4.68e-06 & 339 & 43392 & 20.651 \\
\bottomrule
\end{tabular}

\par\smallskip\footnotesize One frozen jet is used per dimension, at both tolerances. At $d=8,16$, these are gradients of retained learned critics at recorded states; displayed gaps include interval neural-derivative allowances. At larger dimensions they are independently generated jets for the same Hamiltonian family. The experiment measures certificate work for these instances, not an average-case scaling law or a whole-state neural PDE residual.
\end{table}

The 128-dimensional $10^{-5}$ calculation uses 339 scalar boxes and 43,392 scalar convex subproblems. It replaces neither raw actor training nor the preceding finite tensor safeguard in retrospect. At the learned eight- and sixteen-dimensional jets, interval forward differentiation of the stored critic weights bounds the $\ell^1$ gradient discrepancies below $4.584\times10^{-17}$ and $7.877\times10^{-17}$, respectively. These allowances are propagated into the action gap. Theorem~\ref{r9:bridge} identifies the additional state-cover and residual quantities that would be needed for a whole-policy guarantee.

\section{An absolute bound and a frozen external comparison}\label{r9:external}
\subsection{A rigorous lower bound for the unchanged nonconvex family}
Let $r_0=1+.1\sin^2(1)$, $\sigma^2=.32$, and $\lambda_d=\sigma^2r_0/(2d)$. Since $\sin^2a\geq\sin^2(1)a^2$ for $|a|\leq1$, replacing the running cost by $r_0\|a\|^2/d-.3$ and allowing unbounded controls gives a relaxation. The original terminal function $q_d$ is retained.
\begin{proposition}[Scalar-integral absolute lower bound]\label{r9:lower}
Let $Z$ be normal with mean $2\sqrt d/(r_0+4)$ and variance $\sigma^2r_0/(r_0+4)$. Then
\begin{equation}
 V_d^*\geq
 \frac{.25r_0}{r_0+4}+\frac{\sigma^2r_0}{4}\log(1+4/r_0)-.3
 -\lambda_d\log\E\exp\!\left[-\frac{.2}{\lambda_d}\cos Z\right].\label{r9:lowerformula}
\end{equation}
An outward enclosure of the positive expectation yields a rigorous numerical lower bound for the original objective \eqref{r9:nonconvex}.
\end{proposition}
The relaxation is an upper/lower-bound device, not a new training problem with a known manufactured value. The terminal cosine, which couples all state coordinates, remains in the scalar expectation. Its tails are bounded analytically and its interior is integrated by interval rectangles.

\begin{table}[t]
\centering\scriptsize
\caption{Absolute lower bounds for the original nonconvex control problem}\label{r9:tab:absolute}
\begin{tabular}{r r r r r}
\toprule
$d$ & Quadrature cells & Relaxed-value enclosure & Original-value lower & Seconds \\
\midrule
8 & 1024 & $[-0.08599082,-0.08384626]$ & -0.08599082 & 0.260 \\
8 & 4096 & $[-0.08518701,-0.08465086]$ & -0.08518701 & 1.081 \\
8 & 16384 & $[-0.08498598,-0.08485194]$ & -0.08498598 & 4.506 \\
16 & 1024 & $[-0.20520651,-0.20360335]$ & -0.20520651 & 0.271 \\
16 & 4096 & $[-0.20460665,-0.20420585]$ & -0.20460665 & 1.067 \\
16 & 16384 & $[-0.20445644,-0.20435623]$ & -0.20445644 & 4.212 \\
32 & 1024 & $[-0.28757809,-0.28684340]$ & -0.28757809 & 0.245 \\
32 & 4096 & $[-0.28730394,-0.28712026]$ & -0.28730394 & 0.988 \\
32 & 16384 & $[-0.28723515,-0.28718923]$ & -0.28723515 & 4.172 \\
\bottomrule
\end{tabular}

\par\smallskip\footnotesize A row encloses the relaxed value. Its lower endpoint is also a lower bound for the original bounded-control value. The width shown by successive enclosures is quadrature uncertainty; it does not include or estimate the structural relaxation gap.
\end{table}

\subsection{Holdout design and exact deployment interpretation}
The comparison retains the original actor and critic architectures: two tanh hidden layers with 48 units each, a bounded actor output, and the same terminal-value embedding. NBO uses separated evaluation and improvement with Adam, three critic steps per actor step, 256 training states, and four successors. SOC-MartNet uses its pinned author solver, RMSprop, the 600-output test network, and the inherited $J=2$, $K=1$ multiplier setup. These different optimization systems are disclosed rather than described as identical implementations.

The committed protocol freezes dimensions $8,16,32$ and training budgets of 10 and 30 seconds. At each dimension, both methods receive the same six validation trials: two seeds and three learning-rate multipliers. Selection minimizes mean independent validation cost, with a fixed tie rule. Twelve different holdout seeds, 720--731, are then evaluated at each dimension--budget pair. Within each pair, actor and critic starting weights are identical and their hashes are checked. Framework warmup precedes both clocks, and method order is seed-randomized. Training-only and setup-plus-training clocks are retained. The baseline $a_i(t,x)=\operatorname{proj}_{[-1,1]}[-2(x_i-.5)/\{1+4(1-t)\}]$ is specified independently from the associated quadratic problem and is evaluated on the unchanged nonquadratic objective.

Each trained network is deployed as a sample-and-hold control on 128 dates and evaluated with 4,096 independent, paired Gaussian paths. In \eqref{r9:nonconvex}, drift under a held action and volatility are constant on a slab; the Gaussian transition and running-cost integral are therefore exact in real arithmetic for \emph{that deployed continuous-time policy}. A 64-date evaluation uses aggregated common increments and changes the deployment policy. Its difference from the 128-date result is not called an Euler-error bound. Floating-point network evaluation and Monte Carlo uncertainty remain distinct from the outward deterministic lower bound in Table~\ref{r9:tab:absolute}.


\begin{table}[t]\centering\scriptsize
\caption{Frozen holdout comparison: twelve seed pairs per panel}\label{r9:tab:external}
\begin{tabular}{r r r r r r r}
\toprule
$d$ & Budget & NBO cost & SOC cost & LQ cost & Median difference & Adjusted $p$ \\
\midrule
8 & 10 & 0.5812 & 1.0842 & 0.5121 & -0.4783 & 0.001465 \\
8 & 30 & 0.5785 & 1.0557 & 0.5121 & -0.4348 & 0.001465 \\
16 & 10 & 0.4663 & 0.7893 & 0.1740 & -0.1672 & 1 \\
16 & 30 & 0.4630 & 0.8456 & 0.1740 & -0.2878 & 0.438 \\
32 & 10 & 0.3813 & 1.2011 & 0.0854 & -0.8691 & 0.001465 \\
32 & 30 & 0.3799 & 1.0218 & 0.0854 & -0.6962 & 0.001465 \\
\bottomrule
\end{tabular}

\par\smallskip\footnotesize Smaller realized cost is better. The median difference is NBO minus SOC-MartNet. The last column is a one-sided exact sign-test $p$-value adjusted across the six predeclared dimension--budget panels. The sign test concerns the joint training-and-audit pipeline under independent seed draws, not an outward numerical error guarantee.
\end{table}
\begin{table}[t]\centering\scriptsize
\caption{Absolute reference excess and actual training clocks}\label{r9:tab:externalabsolute}
\begin{tabular}{r r r r r r r}
\toprule
$d$ & Budget & Lower bound & NBO excess & SOC excess & NBO sec. & SOC sec. \\
\midrule
8 & 10 & -0.08499 & 0.6662 & 1.1692 & 10.002 & 10.006 \\
8 & 30 & -0.08499 & 0.6635 & 1.1407 & 30.004 & 30.008 \\
16 & 10 & -0.20446 & 0.6707 & 0.9938 & 10.002 & 10.005 \\
16 & 30 & -0.20446 & 0.6674 & 1.0500 & 30.004 & 30.006 \\
32 & 10 & -0.28724 & 0.6686 & 1.4883 & 10.003 & 10.007 \\
32 & 30 & -0.28724 & 0.6672 & 1.3090 & 30.003 & 30.007 \\
\bottomrule
\end{tabular}

\par\smallskip\footnotesize Excess is the estimated mean policy cost minus the rigorous original-value lower bound. It provides an absolute reference axis but is not itself a deterministic regret enclosure: sampling error and floating-point rollout evaluation are not outward certified. Training clocks include completed-update overshoot and exclude separately recorded warmup and setup.
\end{table}


The holdout results favor NBO over the author implementation in mean cost in all six panels, but their interpretation is not uniform. In dimensions 8 and 32, NBO has lower paired cost for all twelve seeds at both budgets; each multiplicity-adjusted sign-test value is $.001465$. In dimension 16 the adjusted values are 1 and $.438$, so these panels do not establish a robust seed-level ordering under the declared test. Tripling the budget gives only small improvements in NBO mean cost, from $.5812$ to $.5785$, $.4663$ to $.4630$, and $.3813$ to $.3799$ respectively. Most importantly, the independently specified quadratic-feedback baseline has lower mean cost than both learned methods in every panel. Its mean costs are $.5121$, $.1740$, and $.08544$ in dimensions 8, 16, and 32. These retained results support the stated NBO--SOC-MartNet comparison at the frozen budgets, not a claim that NBO is the best available control rule or that additional training has already reached an accurate optimum.

The raw records retain every seed difference, the median and interquartile range, approximate Student intervals, costs for all three rules, path standard errors, training curves, optimizer clocks, and checkpoints. Approximate Student intervals are descriptive; the exact sign calculation avoids assuming Gaussian seed-level differences. The source protocol predates holdout execution, while the benchmark family and adapters themselves were developed in earlier revisions. Holdout seeds do not turn one problem family into evidence of general superiority across economic environments. Two budget points are not a matched-accuracy frontier, and a finite absolute lower bound need not be tight enough to rank methods at a small prescribed regret.

\section{Method-specific tuning and an independent replication}\label{r11:robustness}
The preceding holdout answers a precisely frozen comparison, but its scalar learning-rate search leaves an important design question: does the ordering survive a more consequential method-specific tuning effort? We answer this question with a new experiment, rather than choosing new hyperparameters on the already observed holdout or suppressing the earlier results. The objective, diffusion, action box, horizon, equal actor/critic initializations, author solver commit, and 128-date deployed policy are unchanged. The new protocol fixes tuning seeds 1600 and 1601 and holdout seeds 1820--1831 before execution. None of these is an R9 holdout seed.

\subsection{Prospectively specified equal tuning opportunities}
For each dimension and method, twelve initial configurations and two prospectively specified extension configurations are each trained for five nominal seconds on two tuning seeds. Thus each method receives 140 nominal training seconds per dimension. Completed-update overshoot, initialization, validation, and holdout clocks are recorded separately. The methods receive equal nominal training opportunity, not equal numbers of gradient operations or an assertion that their entire hyperparameter spaces have been optimized.

The shared learning-rate multiplier set expands from $\{1/3,1,3\}$ to $\{1/3,1,3,9,27\}$. NBO additionally varies evaluation updates per actor improvement and transition particles. SOC-MartNet additionally varies the adversary learning rate, the number of sinusoidal test outputs, the solver's $J/K$ parameters, and the initial and updating multiplier parameters. The complete twelve-configuration lists appear in the frozen executable protocol; unchanged settings and unsuccessful candidates are not omitted. The pinned author training implementation is unmodified. Method order is randomized within tuning seed and candidate, and actor and critic initial hashes must agree within each holdout pair.

After the first twelve configurations, selection minimizes the mean validation cost over the two tuning seeds, with an index-based tie rule. A selected multiplier at the current high or low boundary is extended by a factor of three in that direction. If the multiplier is interior, the additional slot varies transition-particle count for NBO or adversary rate for SOC. Selection is repeated before the second extension. These rules are specified before execution and do not inspect the new holdout. Failed tuning trials receive infinite selection score while their raw records are preserved. Validation uses 1,024 common Brownian paths per tuning seed, and holdout evaluation uses 4,096 new common paths per seed. All recorded outcomes in this execution are finite, and all 168 tuning trials and 36 holdout pairs are retained.

\begin{table}[t]\centering\small
\caption{Selected configurations and actual tuning-training opportunity}\label{r11:tab:tuning}
\begin{tabular}{rllrr}
\toprule
$d$ & NBO selected & SOC selected & NBO sec. & SOC sec. \\
\midrule
8 & 3 ($m=9$) & 11 ($m=3$) & 140.15 & 140.29 \\
16 & 12 ($m=3$) & 2 ($m=3$) & 140.09 & 140.18 \\
32 & 1 ($m=1$) & 3 ($m=9$) & 140.17 & 140.31 \\
\bottomrule
\end{tabular}

\par\smallskip\footnotesize Configuration indices are zero based and bind to the complete frozen lists and extension records. $m$ is the actor/critic learning-rate multiplier. Actual seconds sum all 28 tuning fits per method and dimension, including completed-update overshoot but excluding separately retained validation and setup.
\end{table}

All six selected learning-rate multipliers are interior to the enlarged tested multiplier set. At $d=8$, NBO selects $m=9$, while SOC selects $m=3$ with initial multiplier and multiplier step both 30. At $d=16$, NBO selects six evaluation steps and sixteen transition particles with $m=3$; SOC selects its default method-specific parameters with $m=3$. At $d=32$, NBO selects $m=1$ and SOC selects $m=9$, with the respective default remaining parameters. The large learning-rate boundary identified in the earlier report is therefore tested beyond its old limit. Interior selection on this finite set does not establish global tuning optimality or exclude improvements along an unsearched interaction. The experiment tests a finite, disclosed method-specific tuning budget.

\subsection{Independent holdout and protocol sensitivity}
The selected configurations are fixed before their twelve new holdout pairs are run at a ten-second nominal training budget. The primary test is again the one-sided exact sign test, now with Bonferroni correction across the \emph{three new} predeclared panels. This does not change the six-panel correction of the earlier experiment. Because the historical evidence motivated the new experiment, the new study is a prospective post-review replication, not a retroactive preregistration of the original benchmark design.

\begin{table}[t]\centering\scriptsize
\caption{New holdout after equal-budget method-specific tuning}\label{r11:tab:robustness}
\begin{tabular}{rrrrrrrr}
\toprule
$d$ & NBO & SOC & LQ & Mean diff. & Median diff. & Wins & Adj. $p$ \\
\midrule
8 & 0.5542 & 0.7928 & 0.5125 & -0.2386 & -0.1967 & 11/12 & 0.00952 \\
16 & 0.4547 & 0.8067 & 0.1741 & -0.3520 & -0.1212 & 8/12 & 0.58154 \\
32 & 0.3824 & 0.4586 & 0.0852 & -0.0762 & -0.0427 & 7/12 & 1.00000 \\
\bottomrule
\end{tabular}

\par\smallskip\footnotesize Smaller realized cost is better. Differences are NBO minus SOC, using paired audit paths and twelve independently seeded training-and-audit pipelines. The last column applies the prospectively fixed three-panel sign-test correction. Costs are sampled policy evaluations, not outward deterministic certificates. The LQ column retains the independently specified clipped feedback under the unchanged nonconvex objective.
\end{table}

The NBO mean differences are $-.23863$, $-.35199$, and $-.07623$ in dimensions 8, 16, and 32. The corresponding paired win counts are 11, 8, and 7 out of twelve, with adjusted sign-test values $.00952$, $.58154$, and 1. Only the eight-dimensional panel establishes the declared robust seed-level ordering at level $.05$. In dimension 16, four seeds again favor SOC. In dimension 32, the expanded tuning yields a substantially smaller mean difference and no significant directional ordering, although the earlier frozen experiment favored NBO in all twelve seeds. The improved SOC result is consistent with its selected $m=9$, which the old scalar grid excluded; this association is not an isolated causal estimate of that hyperparameter because tuning, validation seeds, and holdout seeds all differ.

The independently specified clipped-feedback policy has lower mean realized cost than either trained pipeline in each new panel, as it did in the previous experiment. Its favorable performance does not make its cost a lower bound for a minimization problem, and it does not prove global optimality. These outcomes distinguish three questions: whether a configured NBO pipeline improves on a configured external solver, whether the ordering is robust to a changed tuning protocol, and whether either neural policy has small absolute error. The new experiment supplies direct evidence for the second question and retains all evidence relevant to the first; Section~\ref{r11:accuracy} supplies a separate exact-reference experiment for the third. Favorable absolute-accuracy results in another model are not used to erase the sixteen- or thirty-two-dimensional sign-test outcomes here.

The evidence pipeline collects with an unconditional completion policy. Every planned cell has explicit status before dependency installation; successful outputs, failed outputs, and missing-cell metadata are versioned before aggregate validity is checked. Separate immutable commits identify the frozen study source and completed raw results. No branch name alone is used as a review target, and the 36 new pairs neither replace nor pool selectively with the 72 earlier pairs.


\section{From implemented actor errors to absolute-accuracy frontiers}\label{r11:accuracy}
The original stopped economy establishes economically resolved optimality bounds. A complementary experiment asks a different question: how much computation is required to attain a common absolute accuracy when the continuous optimum is independently known? This experiment changes neither the economy in Section~\ref{r9:economy} nor the coupled nonconvex comparison. It adds reference-solvable mechanisms with different drift, noise, and control-cost conditioning. An exact reference is not replaced by the best realized cost of a competing policy.

Consider independent coordinates with
\begin{equation}\label{r11:lq}\begin{aligned}
 dX^j_t&=(a_jX^j_t+b_ju^j_t)dt+\sigma_jdW^j_t,\\
 J(u)&=\frac1d\E\left[\int_0^1\sum_{j=1}^d(q_j(X^j_t)^2+r_j(u^j_t)^2)dt+\sum_{j=1}^dw_j(X^j_1)^2\right],
\end{aligned}\end{equation}
where $q_j\geq0$, $r_j,w_j>0$, and $b_j>0$. Controls are unbounded, adapted, and square integrable. The action constraints and nonconvex running cost of the external comparison are \emph{not} silently imposed on this linear-quadratic problem. Rather, the latter supplies an independently solvable accuracy experiment. The implemented actor observes $X_{i/N}$ and holds $u_t=-K_iX_{i/N}$ until the next decision date. It is a linear output map with time gates, not a continuous-feedback deployment of $-K_iX_t$ and not the width-48 tanh actor used in the external comparison.

\begin{proposition}[Two independent error decompositions]\label{r11:decomposition}
Let $p_j(1)=w_j$ and $-\dot p_j=q_j+2a_jp_j-b_j^2p_j^2/r_j$. The continuous optimum is
\begin{equation}\label{r11:reference}
 V^*=\frac1d\sum_j\left[p_j(0)(X^j_0)^2+\sigma_j^2\int_0^1p_j(t)dt\right],
\end{equation}
and every admissible control satisfies
\begin{equation}\label{r11:continuousdefect}
 J(u)-V^*=\frac1d\sum_j\E\int_0^1r_j\left(u^j_t+\frac{b_jp_j(t)}{r_j}X^j_t\right)^2dt.
\end{equation}
For decision interval $h=1/N$, let $J_h^*$ be the optimum over held-action rules, with Riccati coefficients $P_{ij}$, optimal gains $K^*_{ij}$, and positive action curvatures $D_{ij}$ specified in Appendix~\ref{r11:proof}. For any implemented held-action rule, including a nonlinear neural rule,
\begin{equation}\label{r11:twodefect}
 J(u)-V^*=(J_h^*-V^*)+
 \frac1d\sum_{i,j}D_{ij}\E\left[(u^j_i+K^*_{ij}X^j_{ih})^2\right].
\end{equation}
Consequently, output-error bounds $\E[(u_i^j+K^*_{ij}X^j_{ih})^2]\leq\varepsilon_{ij}^2$ give a quantitative optimization allowance $d^{-1}\sum_{i,j}D_{ij}\varepsilon_{ij}^2$, in addition to the separately evaluated deployment error $J_h^*-V^*$. No convergence assumption on Adam or on shared neural parameters is used.
\end{proposition}

The proposition is a specialization of quadratic verification and Bellman completion of squares, not a new general Riccati theorem. Its role is to connect an \emph{implemented} policy error to an absolute economic objective. The finite-horizon quadratic Bellman construction requires exactly $N$ backward improvements per coordinate. In this output class, evaluation and improvement are closed operations, and the backward recursion solves the held-action problem rather than merely decreasing a stochastic training loss. The alternative Adam implementation optimizes the same gain outputs using an independently coded exact-moment objective. Its errors are assessed after the recorded updates; a universal stochastic-optimization convergence theorem is not inferred.

\subsection{Independent reference and deployment verification}
Write $c_j=b_j^2/r_j$, $\lambda_j^2=a_j^2+c_jq_j$, and $s=1-t$. Define
\begin{align*}
 A_j(s)&=\cosh(\lambda_js)+\frac{\sinh(\lambda_js)}{\lambda_j}(-a_j+c_jw_j),\\
 B_j(s)&=w_j\cosh(\lambda_js)+\frac{\sinh(\lambda_js)}{\lambda_j}(q_j+a_jw_j).
\end{align*}
The continuous Riccati solution is $p_j(t)=B_j(1-t)/A_j(1-t)$, with the continuous limit when $\lambda_j=0$. Its integral equals $(\log A_j(1)+a_j)/c_j$. Thus the reference needs no time-stepping solver or simulated paths. Sixty-decimal interval arithmetic encloses the exponentials, logarithms, and divisions. Every coefficient in the certificate is the exact binary64 number stored in the instance file, and every gain is the exact stored policy output. Exported interval endpoints are enlarged by one binary64 unit in the outward direction.

A different calculation evaluates the original continuous cost of each held-action policy. Conditional on the state at a decision date,
\[
 X^j_{ih+s}=\left(e^{a_js}-b_jK_{ij}\int_0^s e^{a_jv}dv\right)X^j_{ih}
       +\sigma_j\int_0^s e^{a_j(s-v)}dW^j_{ih+v}.
\]
Exact second-moment propagation and the integrated running cost therefore determine $J(u)$. Appendix~\ref{r11:proof} gives all coefficients and the verification proof. This is not an Euler approximation, and a difference between two meshes is not used as an error estimate. In particular, the reported interval $[\underline J-\overline V,\overline J-\underline V]$ encloses absolute regret in \eqref{r11:lq}, including deployment error. Stored policies, model coefficients, source hashes, and the two independent evaluators are supplied in the replication package.

\subsection{Reference-solvable mechanisms and executed accuracy}
The prospective protocol specifies three cases. Tracking has $d=8$, $a=0$, $b=2$, $q=0$, $r=w=1$, initial centered state $-.5$, and $\sigma=.5656854249492381$. The mean-reverting case has $d=8$, $a=-1$, $b=q=w=1$, $r=.25$, $\sigma=.4$, and initial state one. The ill-conditioned case has $d=16$ and linearly spaced $a\in[-1,.5]$, $b\in[.5,2]$, $q\in[.1,2]$, $\sigma\in[.15,.65]$, $w\in[.5,2]$, and initial states in $[-1,1]$, with geometrically spaced $r\in[.1,4]$. The last case contains both stable and unstable uncontrolled coordinates. Exact stored coefficients, rather than rounded endpoints in this description, define each computation.

The quadratic Bellman constructor is evaluated at $N=4,8,16,32,64,128$. Adam uses $N=32$, initially zero gains, learning rate $.05$, standard moment coefficients $(.9,.999)$, and checkpoints at 0, 25, 100, and 400 updates. The deterministic optimizer configuration is recorded with the executed instances; no tuning search is performed for this reference suite. No holdout claim is attached to these analytical test cases. All checkpoints are retained, including those that miss a stated accuracy threshold.

\begin{table}[t]\centering\small
\caption{Absolute continuous-model regret: outward upper endpoints}\label{r11:tab:accuracy}
\begin{tabular}{lrrrr}
\toprule
Case & B: $N=16$ & B: $N=32$ & B: $N=128$ & Adam (400) \\
\midrule
Tracking & 0.00839748 & 0.00409984 & 0.00100625 & 0.00409984 \\
Mean-reverting & 0.00429055 & 0.00211098 & 0.00052088 & 0.00211101 \\
Ill-conditioned & 0.01421242 & 0.00699643 & 0.00172851 & 0.00699643 \\
\bottomrule
\end{tabular}

\par\smallskip\footnotesize Each entry is rounded upward from an interval certificate of the original continuous cost of the stored actor. Columns B denote quadratic Bellman construction at the stated decision count. Adam uses 32 decisions and 400 updates. These are structured-policy accuracy results, not a substitution for the external neural comparison.
\end{table}

Table~\ref{r11:tab:accuracy} shows that refinement reduces deployment error in every case. At 128 decisions, all three absolute regret upper bounds are below $.003$. At 32 decisions, all three are below $.01$. Adam approaches the same 32-decision optimum: increasing its updates cannot eliminate the deployment floor at that fixed decision count. This distinction is observed directly, not inferred from a nearly flat training curve. The $.001$ target is attained on the planned mesh sequence for the mean-reverting case, but not for tracking or the ill-conditioned case; these outcomes remain in the package rather than being removed from the frontier.

\begin{table}[t]\centering\small
\caption{First certified attainment on the recorded computation paths}\label{r11:tab:frontier}
\begin{tabular}{llrrr}
\toprule
Case & Method & $.01$ (sec.) & $.003$ (sec.) & $.001$ (sec.) \\
\midrule
Tracking & B & 0.023 & 0.086 & --- \\
Tracking & A & 0.796 & --- & --- \\
Mean-reverting & B & 0.017 & 0.059 & 0.193 \\
Mean-reverting & A & 0.820 & 3.046 & --- \\
Ill-conditioned & B & 0.099 & 0.374 & --- \\
Ill-conditioned & A & 1.019 & --- & --- \\
\bottomrule
\end{tabular}

\par\smallskip\footnotesize Time includes all earlier generation and certificate checks on the indicated path, but excludes dependency installation and one-time optimizer construction. B is the quadratic Bellman constructor; A is Adam on gain outputs. A dash means the target was not attained among the predeclared checkpoints. Timing is a single-machine descriptive measurement, not a confidence interval or a hardware-independent complexity comparison.
\end{table}

The accuracy axis is the same original-cost regret for both policy generators within each case. The computation axis includes rejected earlier meshes or checkpoints rather than reporting only the final inexpensive verification. The exact quadratic constructor uses analytic model structure that the general tanh and SOC implementations do not exploit; its faster times do not establish a neural-versus-SOC efficiency ranking. Together with the unchanged stopped economy and coupled nonconvex family, the added cases test first-exit verification, adapted adjustment, nonconvex coupling, mean reversion, and differing control-cost conditioning. They do not supply new recursive-utility or game experiments. Those extensions retain their explicit assumptions and complete historical material.


\section{Relation to rigorous control and neural computation}\label{r9:related}
The strongest claims should be compared with methods that establish bounds, not only with methods that train networks. Theorem~\ref{r9:pair} is an application of classical stopped verification to smooth upper and lower witnesses. Its role here is to handle the actual liquidation trace without presupposing a smooth or boundary-regular optimal value. Theorem~\ref{r9:bridge} supplies a computation-facing cover, derivative, and oracle allowance. Neither theorem claims that classical verification was absent from the literature.

Monotone, stable, consistent approximations with comparison provide convergence to viscosity solutions in the framework of \citet{barles1991}. Controlled Markov-chain approximation is developed systematically by \citet{kushner2001}; adaptive sparse-grid methods address high-dimensional dynamic models in \citet{brumm2017}. Such results concern specified approximation schemes and their hypotheses. The finite telescoping budget in the preserved manuscript names the operator discrepancies needed for a continuous statement; an observed mesh difference does not supply them. The original-payoff certificates here instead evaluate continuous generators or exact Gaussian moments directly, with separately bounded analytic and arithmetic remainders. They do not assert a convergence rate for the inherited killed-Euler scheme.

Primal--dual numerical control methods also pair feasible decisions with bounds. \citet{brown2010} relax nonanticipativity with penalties and establish weak and strong duality for stochastic dynamic programs. \citet{brown2011} use related bounds in portfolio optimization with transaction costs. Our restricted and flexible constructions are market-deflator witnesses with explicit stopped localization; they are not claims that policy-plus-upper-bound analysis is new. Theorem~\ref{r10:dual} adds a verified affine-preference source and a conditional-coefficient variance allowance. Its distinction from an information-relaxation bound with an uncomputed penalty is that every source integral, variance constant, and face allowance is evaluated here. \citet{lasserre2008} develop occupation-measure and moment relaxations for polynomial nonlinear optimal control under their stated assumptions. We do not invoke their deterministic convergence theorem for the present logarithmic, stochastic, controlled-exit problem. The explicit source bounds, interval polynomial conversions, Gaussian tails, and trace inequalities are verified here instead.

Theorem~\ref{r10:dual}, Proposition~\ref{r9:upper}, and Theorem~\ref{r9:restricted} identify the current and preceding model-specific constructions behind the economic bounds. The affine-preference theorem retains costly adjustment through an adapted conditional coefficient, bounds its randomness quantitatively, and closes the declared central-state target. The fitted generator proposition couples a polynomial witness to a small localization potential; the restricted theorem cancels every portfolio term before computing a bounded-source expectation. Theorem~\ref{r9:actionthm} exploits rank-one action coupling and strong convexity of the remaining scalar problems. It is a structured convex-duality and branch-and-bound construction, not a generic replacement for global nonconvex optimization. Proposition~\ref{r9:lower} preserves the nonquadratic terminal payoff in a computable relaxation of the exact comparison problem.

Neural policy iteration and neural PDE solvers are important computational antecedents. \citet{jacka2015} analyze continuous-time policy improvement. \citet{kim2026explore} separate iteration, policy-network, and residual errors in entropy-regularized control; \citet{kim2026interior} study interior bounds and the mismatch between collocation and closed-loop performance. \citet{lee2025} combine policy iteration and operator learning. \citet{cai2025published} impose martingale conditions without an explicit Hamiltonian infimum. NBO's raw training comparison is with a pinned implementation of the last method, not a locally recreated approximation described as author code. Our certificates can also audit a policy produced by another solver. This portability is a property of independent verification, while the training results remain evidence about the specifically recorded NBO and SOC-MartNet adapters.

Recursive utility, temporal-self equilibria, and games require their own evaluation operators and optimality notions. Their complete equations, assumptions, finite gain bounds, failures, and earlier numerical results are preserved in the supplement. No new large-scale recursive or game accuracy result is inferred from the additive-utility certificates above. The stochastic-trace bias identities and local-optimizer counterexamples are preserved for the same reason: they identify conditions that a neural training loss cannot replace.

\section{Conclusion}
The original continuous preference-adjustment economy now has certified central-state policies at the declared $.01$ regret tolerance for all three cost coefficients. A feasible learned-and-polished policy supplies each lower bound, while an independent affine-preference dual supplies an upper bound over the original continuous adaptive controls. This closes the numerical target without deleting the original stopping contract or replacing it by a manufactured-payoff problem.

The cost-continuation construction extends this guarantee from three isolated prices to every real price between $.5$ and $8$. Its eighteen-policy library has a largest rigorously enclosed regret of $.00999403$. Only a finite set of exact envelope breakpoints is checked, while convexity and fixed-policy cost transfer establish the continuum statement. The policy is selected once from the library for the given price; the guarantee remains a central-initial-state result, not an uncomputed whole-state neural certificate.

The value of access and the welfare cost of more expensive adjustment are resolved at a precision calibrated to the reported effects. At coefficient two the optimal access gain is at least $.04096859$; raising the coefficient to eight lowers optimal welfare by at least $.02318257$. Both guarantees exceed their corresponding interval widths. Optimal adjustment budgets obey the analytical ordering and a valid secant enclosure, but are not equated with the budgets of approximate policies.

The constructive tools also retain their distinct roles: fitted Bernstein witnesses, stopped Gaussian-moment evaluation, market-deflator localization, quantitative output and neural-derivative allowances, and a structured continuous-action oracle. The unchanged coupled comparison problem has an absolute lower reference and a complete frozen holdout record, including unfavorable baseline and sign-test outcomes. Those experiments do not establish a universally best-tuned neural solver or an external matched-accuracy dominance theorem. The added reference experiments give certified continuous-model accuracy, implemented-output error identities, and computation-to-accuracy paths rather than only fixed-budget losses. The expanded tuning study assesses robustness under a second frozen protocol, not by retuning the original holdout. These additional results strengthen the executable theory--algorithm link while leaving each certificate attached to its actual model, policy class, and implementation. All preceding manuscripts, equations, negative diagnostics, and raw outcomes remain preserved for further examination.
